% SOURCE_AUTHORITY: sga5_fr_workpass.tex, lines 2745--2893 (LF-normalized unit).
% SOURCE_SHA256: 791F4EFFC5E02832D5D77ED03518C8156D6F07E4C8238B03545DB93D883FBB28
\subsection*{4.1}
Sean, para $i=1,2$, $X_i$ un $S$-esquema, $X=X_1\times_SX_2$, $p_i:X\to X_i$ las proyecciones, $L_i\in\ob\Dctf(X_i)$ y $K_{X_i}$ el complejo dualizante (1.3). Las flechas canónicas
\begin{equation}
p_i^*\uRHom(L_i,K_{X_i})\overset{\mathbf L}{\otimes}p_i^*L_i
\longrightarrow p_i^*K_{X_i},
\tag{4.1.1}
\end{equation}
compuestas de 2.1.1 y de la flecha de evaluación 2.2.2, proporcionan, tomando el producto tensorial sobre $X$ y componiendo con el isomorfismo de Künneth 1.7.6, un apareamiento
\begin{equation}
(DL_1\overset{\mathbf L}{\otimes}_S L_2)\overset{\mathbf L}{\otimes}(L_1\overset{\mathbf L}{\otimes}_S DL_2)
\longrightarrow K_X,
\tag{4.1.2}
\end{equation}
que puede reescribirse, gracias a 3.1.1, 3.1.2,
\begin{equation}
\uRHom_S(L_1,L_2)\overset{\mathbf L}{\otimes}\uRHom_S(L_2,L_1)
\longrightarrow K_X.
\tag{4.1.3}
\end{equation}
De ello se deduce un apareamiento, denominado producto cup,
\begin{equation}
\langle\ ,\ \rangle:\Hom_S(L_1,L_2)\otimes\Hom_S(L_2,L_1)
\longrightarrow \H^0(X,K_X).
\tag{4.1.4}
\end{equation}
También se puede considerar 4.1.2 como el producto tensorial externo de las flechas de evaluación
\[
DL_i\overset{\mathbf L}{\otimes}L_i\longrightarrow K_{X_i}.
\]
Teniendo en cuenta el isomorfismo de Künneth 2.2.4, se ve así que 4.1.2 es una ``dualidad perfecta'', es decir, identifica, mediante el isomorfismo de adjunción 2.2.1, cada uno de los complejos $DL_1\overset{\mathbf L}{\otimes}_S L_2$, $L_1\overset{\mathbf L}{\otimes}_S DL_2$ con el dual del otro (con valores en $K_X$).

Cuando $X_1=X_2=S$, 4.1.3 es un caso particular de un apareamiento válido para complejos perfectos sobre un topos anillado cualquiera (SGA 6 I 7.8), y, con las notaciones de (SGA 6 I 8.3), se tiene
\begin{equation}
\langle u,v\rangle=\mathrm{Tr}(vu)=\mathrm{Tr}(uv).
\tag{4.1.5}
\end{equation}

\subsection*{4.2}
Sean $c:C\to X$, $d:D\to X$ y $e=c\times_Xd:E=C\times_XD\to X$. Para $P,Q\in\ob\Dctf(X)$, se dispone de una flecha canónica (que no es un isomorfismo en general)
\begin{equation}
c^!P\overset{\mathbf L}{\otimes}_X d^!Q
\longrightarrow e^!(P\overset{\mathbf L}{\otimes}Q),
\tag{4.2.1}
\end{equation}
definida, del mismo modo que 1.7.3, como adjunta de la flecha compuesta
\[
e_!(c^!P\overset{\mathbf L}{\otimes}_X d^!Q)
\xrightarrow[(1)]{\sim}
 c_!c^!P\overset{\mathbf L}{\otimes}d_!d^!Q
\xrightarrow{(2)} P\overset{\mathbf L}{\otimes}Q,
\]
donde (1) es el isomorfismo de Künneth (SGA 4 XVII 5.4.2.2) y (2) el producto tensorial de las flechas de adjunción. Aplicando $e_!$ a 4.2.1 y componiendo con el isomorfismo de Künneth que acabamos de citar, se obtiene una flecha canónica
\begin{equation}
c_!c^!P\overset{\mathbf L}{\otimes}d_!d^!Q
\longrightarrow e_!e^!(P\overset{\mathbf L}{\otimes}Q).
\tag{4.2.2}
\end{equation}
Del mismo modo, componiendo la flecha de Künneth
\[
c_*c^!P\overset{\mathbf L}{\otimes}d_*d^!Q
\longrightarrow e_*(c^!P\overset{\mathbf L}{\otimes}_X d^!Q)
\]
con la flecha deducida de 4.2.1 aplicando $e_*$, se obtiene una flecha canónica
\begin{equation}
c_*c^!P\overset{\mathbf L}{\otimes}d_*d^!Q
\longrightarrow e_*e^!(P\overset{\mathbf L}{\otimes}Q).
\tag{4.2.3}
\end{equation}
Para $P=DL_1\overset{\mathbf L}{\otimes}_S L_2$, $Q=L_1\overset{\mathbf L}{\otimes}_S DL_2$, de 4.2.2, 4.2.3 se deducen, por composición con 4.1.2 y el isomorfismo $e^!K_X=K_E$, y con las identificaciones 3.1.1, 3.2.1, los apareamientos
\begin{equation}
c_!\uRHom(c_1^*L_1,c_2^!L_2)\overset{\mathbf L}{\otimes}
 d_!\uRHom(d_2^*L_2,d_1^!L_1)
\longrightarrow e_!K_E,
\tag{4.2.4}
\end{equation}
\begin{equation}
c_*\uRHom(c_1^*L_1,c_2^!L_2)\overset{\mathbf L}{\otimes}
 d_*\uRHom(d_2^*L_2,d_1^!L_1)
\longrightarrow e_*K_E,
\tag{4.2.4'}
\end{equation}
\begin{equation}
\langle\ ,\ \rangle:
\Hom(c_1^*L_1,c_2^!L_2)\otimes\Hom(d_2^*L_2,d_1^!L_1)
\longrightarrow \H^0(E,K_E).
\tag{4.2.5}
\end{equation}
Estos apareamientos son compatibles con la localización étale: dado un cuadrado conmutativo (no necesariamente cartesiano)
\[
\begin{tikzcd}[column sep=2.5em,row sep=1.5em]
& E' \arrow[dl] \arrow[dr] & && & E \arrow[dl] \arrow[dr] & \\
C' \arrow[dr] && D' \arrow[dl] & \text{ étale sobre } & C \arrow[dr] && D \arrow[dl] \\
& X &&&& X &
\end{tikzcd}
\]
si denotamos por $c'$, $d'$, $e'$ las compuestas $C'\to C\to X$, $D'\to D\to X$, $E'\to E\to X$, se tiene un cuadrado conmutativo.


\begin{equation}
\tag{4.2.6}
\begin{tikzcd}[column sep=large,row sep=large]
\Hom(c_1^*L_1,c_2^!L_2)\otimes\Hom(d_2^*L_2,d_1^!L_1)
  \arrow[r,"4.2.5"] \arrow[d] &
\H^0(E,K_E)\arrow[d] \\
\Hom(c_1'{}^*L_1,c_2'{}^!L_2)\otimes\Hom(d_2'{}^*L_2,d_1'{}^!L_1)
  \arrow[r] &
\H^0(E',K_{E'})
\end{tikzcd}
\end{equation}
donde las flechas verticales son las restricciones, y la flecha horizontal inferior es la compuesta de la flecha análoga a 4.2.5 y de la restricción
\[
\H^0(E'',K_{E''})\longrightarrow \H^0(E',K_{E'}),
\]
con $E''=C'\times_{X}D'$ (se tienen cuadrados conmutativos análogos con flechas del tipo 4.2.4, 4.2.4', cuya escritura dejamos al lector). Si $Z$ es étale sobre $E$, para
\[
u\in\Hom(c_1^*L_1,c_2^!L_2),\qquad
v\in\Hom(d_2^*L_2,d_1^!L_1),
\]
escribiremos
\begin{equation}
\tag{4.2.7}
\langle u,v\rangle_Z\in \H^0(Z,K_Z)
\end{equation}
para la imagen de $\langle u,v\rangle$ en $\H^0(Z,K_Z)$ por restricción (en la práctica, $Z$ será una componente conexa de $E$). Si $z$ es un punto cerrado aislado de $E$, $\bar z$ un punto geométrico sobre $z$, de 4.2.6 se deduce que
\[
\langle u,v\rangle_z\in \H^0(z,K_z)=A
\]
depende únicamente de la situación obtenida mediante localización estricta en las imágenes de $\bar z$ en $C,D,X$.

\begin{statement}{Ejemplo 4.3}
Supongamos que $X_i$ es liso, de dimensión relativa pura $r_i$, y que $c$ y $d$ son inmersiones cerradas, con $c_2$ (resp. $d_1$) finito. Para $L_1=A_{X_1}$, $L_2=A_{X_2}$, se tiene, por 3.1.3,
\[
DL_1\otimes_S^{\mathbf L}L_2=A_X(r_1)[2r_1],\qquad
L_1\otimes_S^{\mathbf L}DL_2=A_X(r_2)[2r_2],
\]
y 4.1.2 se identifica con el producto
\[
A_X(r_1)[2r_1]\otimes^{\mathbf L}A_X(r_2)[2r_2]\longrightarrow A_X(r)[2r],
\]
donde $r=r_1+r_2$, y por tanto, teniendo en cuenta 3.2.1, 4.2.5 se identifica con el producto cup habitual
\[
\H_C^{2r_1}(X,A(r_1))\otimes \H_D^{2r_2}(X,A(r_2))
\longrightarrow
\H_{C\cap D}^{2r}(X,A(r)).
\]
Tomemos como correspondencia cohomológica de $X_1$ a $X_2$ (resp. de $X_2$ a $X_1$) la clase de $c$ (resp. de $d$) (3.2 a)). De (SGA $4^{1/2}$ Ciclo 2.3.8) se deduce que, en un punto aislado $z$ de $C\cap D$, $\langle \cl(c),\cl(d)\rangle_z$ no es sino la imagen en $A$ de la multiplicidad de intersección de $C$ y $D$ en $z$.
\end{statement}

He aquí ahora el resultado principal de la exposición.
